Throughout, the clauses (ii) to (v) cited are those of Theorem \claimref{synthesis.size-of-an-element-of-the-switch-tower}.

\emph{(o).} The tests of one switch satisfy $\mathrm{on}_{i}^{2} = v_{i}(w_{i}v_{i})w_{i} = v_{i}(1 - v_{i}w_{i})w_{i} = \mathrm{on}_{i}$, likewise $\mathrm{off}_{i}^{2} = \mathrm{off}_{i}$, and $\mathrm{on}_{i}\,\mathrm{off}_{i} = v_{i}w_{i}^{2}v_{i} = 0 = \mathrm{off}_{i}\,\mathrm{on}_{i}$, $\mathrm{on}_{i} + \mathrm{off}_{i} = 1$. Tests of different switches are even elements of algebras on disjoint index sets, so they commute (clause (iv)), and a product of commuting idempotents is an idempotent. For $i < j$ the product $g_{i}g_{j}$ contains $\mathrm{on}_{i}\,\mathrm{off}_{i} = 0$. The sizes follow from clause (iv), applied one switch at a time, and $N(\mathrm{on}_{i}) = N(\mathrm{off}_{i}) = \tfrac12$ (clause (v)). Finally $P_{m-1} = P_{m-1}(\mathrm{on}_{m} + \mathrm{off}_{m}) = g_{m} + P_{m}$ and $P_{0} = 1$.

\emph{(i).} By (o), $E_{n}g_{j}$ is $g_{j}$ for $j \in J$, $j \le n$, and $0$ for $j \notin J$, $j \le n$.

\emph{(ii).} If $J$ is finite, the sum of the $g_{j}$ with $j \in J$ is a solution: it is $E_{n}$ for $n$ at least the largest element of $J$, and $E_{n}g_{j} = 0$ for $j > n$ as well. If the complement of $J$ is finite, $1$ minus the sum of the $g_{j}$ with $j \notin J$ is a solution. Now let $J$ and its complement both be infinite, and suppose $x \in T$ is a solution. It lies in some rung $A_{m}$. Choose $j \in J$ and $j' \notin J$, both larger than $m$. Then
\[
g_{j} = P_{m}\,h, \qquad h = \mathrm{off}_{m+1}\cdots\mathrm{off}_{j-1}\,\mathrm{on}_{j},
\]
where $h$ is an even element of the algebra of the switches $m + 1, \dots, j$, of size $2^{-(j - m)}$ by the argument of (o). The element $c = xP_{m} - P_{m}$ lies in $A_{m}$, and $ch = xg_{j} - g_{j} = 0$. By clause (iv), $0 = N(ch) = N(c)\,N(h)$, so $N(c) = 0$ and $c = 0$ (clause (ii)): $xP_{m} = P_{m}$. The same argument with $j'$, applied to $c' = xP_{m}$, gives $xP_{m} = 0$. So $P_{m} = 0$, which contradicts $N(P_{m}) = 2^{-m}$.

\emph{(iii).} $E_{2n} = g_{2} + g_{4} + \dots + g_{2n}$ is a sum of idempotents with pairwise products zero, so by clause (v) and (o), $N(E_{2n}) = 4^{-1} + \dots + 4^{-n} = \tfrac13(1 - 4^{-n})$, and $E_{2n+2} - E_{2n} = g_{2n+2}$ has size $4^{-(n+1)}$. For $x \in A_{m}$, $N(x) = j/2^{m}$ for an integer $j$ (clause (i)), and $|j/2^{m} - \tfrac13| = |3j - 2^{m}|/(3 \cdot 2^{m}) \ge 1/(3 \cdot 2^{m})$, because $3j - 2^{m}$ is a nonzero integer. By subadditivity (clause (ii)), $N(a) \le N(a - b) + N(b)$ and $N(b) \le N(b - a) + N(a)$ with $N(b - a) = N(a - b)$, so $N(E_{2n} - x) \ge |N(E_{2n}) - N(x)| \ge |N(x) - \tfrac13| - |\tfrac13 - N(E_{2n})| \ge \tfrac{1}{3 \cdot 2^{m}} - \tfrac{1}{3 \cdot 4^{n}}$.
